% Fichier complété par tools/complete_missing_corrections.py.
% Source : PPM/PPM-03/75_TrainA380/75_TrainA380.tex
% Statut : rédigé (7 question(s)).
\begin{corrigebox}[Corrigé rédigé]
\CorrigeQuestion{1}
Les surfaces de référence sont choisies parce qu'elles
                réalisent la mise en position fonctionnelle de la pièce dans l'assemblage.
                La référence primaire supprime trois degrés de liberté, la secondaire deux
                et la tertiaire le dernier; elles doivent être étendues, stables,
                accessibles au contrôle et directement liées aux surfaces d'appui réelles.
\CorrigeQuestion{2}
Les surfaces de référence sont choisies parce qu'elles
                réalisent la mise en position fonctionnelle de la pièce dans l'assemblage.
                La référence primaire supprime trois degrés de liberté, la secondaire deux
                et la tertiaire le dernier; elles doivent être étendues, stables,
                accessibles au contrôle et directement liées aux surfaces d'appui réelles.
\CorrigeQuestion{3}
La spécification se décode en identifiant successivement
                l'élément tolérancé, le symbole géométrique, la valeur et la forme de la
                zone de tolérance, les références ordonnées et les modificateurs. Le
                gabarit est le simulateur géométrique associé : deux plans pour une forme
                ou une orientation de surface, un cylindre pour un axe, ou une zone
                cylindrique/sphérique pour une localisation. La pièce est conforme si
                l'élément extrait reste entièrement dans cette zone après mise en référence.
\CorrigeQuestion{4}
La spécification se décode en identifiant successivement
                l'élément tolérancé, le symbole géométrique, la valeur et la forme de la
                zone de tolérance, les références ordonnées et les modificateurs. Le
                gabarit est le simulateur géométrique associé : deux plans pour une forme
                ou une orientation de surface, un cylindre pour un axe, ou une zone
                cylindrique/sphérique pour une localisation. La pièce est conforme si
                l'élément extrait reste entièrement dans cette zone après mise en référence.
\CorrigeQuestion{5}
La spécification se décode en identifiant successivement
                l'élément tolérancé, le symbole géométrique, la valeur et la forme de la
                zone de tolérance, les références ordonnées et les modificateurs. Le
                gabarit est le simulateur géométrique associé : deux plans pour une forme
                ou une orientation de surface, un cylindre pour un axe, ou une zone
                cylindrique/sphérique pour une localisation. La pièce est conforme si
                l'élément extrait reste entièrement dans cette zone après mise en référence.
\CorrigeQuestion{6}
La spécification se décode en identifiant successivement
                l'élément tolérancé, le symbole géométrique, la valeur et la forme de la
                zone de tolérance, les références ordonnées et les modificateurs. Le
                gabarit est le simulateur géométrique associé : deux plans pour une forme
                ou une orientation de surface, un cylindre pour un axe, ou une zone
                cylindrique/sphérique pour une localisation. La pièce est conforme si
                l'élément extrait reste entièrement dans cette zone après mise en référence.
\CorrigeQuestion{7}
La spécification se décode en identifiant successivement
                l'élément tolérancé, le symbole géométrique, la valeur et la forme de la
                zone de tolérance, les références ordonnées et les modificateurs. Le
                gabarit est le simulateur géométrique associé : deux plans pour une forme
                ou une orientation de surface, un cylindre pour un axe, ou une zone
                cylindrique/sphérique pour une localisation. La pièce est conforme si
                l'élément extrait reste entièrement dans cette zone après mise en référence.
\end{corrigebox}
